%%%%%%%%%%%%%%%%%%%%%%%%%%%%%%%%%%%%%%%%
% datoteka diploma-FRI-vzorec.tex
%
% vzorčna datoteka za pisanje diplomskega dela v formatu LaTeX
% na UL Fakulteti za računalništvo in informatiko
%
% na osnovi starejših verzij vkup spravil Franc Solina, maj 2021
% prvo verzijo je leta 2010 pripravil Gašper Fijavž
%
% za upravljanje z literaturo ta vezija uporablja BibLaTeX
%
% svetujemo uporabo Overleaf.com - na tej spletni implementaciji LaTeXa ta vzorec zagotovo pravilno deluje
%

\documentclass[a4paper,12pt,openright]{book}
%\documentclass[a4paper, 12pt, openright, draft]{book}  Nalogo preverite tudi z opcijo draft, ki pokaže, katere vrstice so predolge! Pozor, v draft opciji, se slike ne pokažejo!
 
\usepackage[utf8]{inputenc}   % omogoča uporabo slovenskih črk kodiranih v formatu UTF-8
\usepackage[slovene,english]{babel}    % naloži, med drugim, slovenske delilne vzorce
\usepackage[pdftex]{graphicx}  % omogoča vlaganje slik različnih formatov
\usepackage{fancyhdr}          % poskrbi, na primer, za glave strani
\usepackage{amssymb}           % dodatni matematični simboli
\usepackage{amsmath}           % eqref, npr.
\usepackage[hyphens]{url}
\usepackage{csquotes}
\usepackage[pdftex, colorlinks=true,
						citecolor=black, filecolor=black, 
						linkcolor=black, urlcolor=black,
						pdfproducer={LaTeX}, pdfcreator={LaTeX}]{hyperref}
\usepackage{hyperxmp}
\usepackage{color}
\usepackage{soul}

\usepackage{listings}
\renewcommand{\lstlistingname}{Izvorna koda}
%\lstset{
%  basicstyle=\fontsize{11}{13}\selectfont\ttfamily
%}
\lstdefinestyle{mystyle}{
  commentstyle=\color{codegreen},
  keywordstyle=\color{magenta},
  %numberstyle=\tiny\color{codegray},
  stringstyle=\color{codepurple},
  basicstyle=\ttfamily\footnotesize,
  breakatwhitespace=false,         
  breaklines=true,                 
  captionpos=b,                    
  keepspaces=true,                 
  numbers=left,                    
  numbersep=5pt,                  
  showspaces=false,                
  showstringspaces=false,
  showtabs=false,                  
  tabsize=2
}
\lstset{style=mystyle}

\usepackage[
backend=bibtex,
style=numeric,
sorting=nty,
]{biblatex}


\addbibresource{literatura.bib} %Imports bibliography file


%%%%%%%%%%%%%%%%%%%%%%%%%%%%%%%%%%%%%%%%
%	DIPLOMA INFO
%%%%%%%%%%%%%%%%%%%%%%%%%%%%%%%%%%%%%%%%
\newcommand{\ttitle}{Inženiring vzporednih algoritmov}
\newcommand{\ttitleEn}{Parallel Algorithm Engineering}
\newcommand{\tsubject}{\ttitle}
\newcommand{\tsubjectEn}{\ttitleEn}
\newcommand{\tauthor}{Luka Bajić}
\newcommand{\tkeywords}{vzporedni algoritmi, večnitnost, Levenshteinova razdalja, LCS}
\newcommand{\tkeywordsEn}{computer, computer, computer}

%%%%%%%%%%%%%%%%%%%%%%%%%%%%%%%%%%%%%%%%
%	HYPERREF SETUP
%%%%%%%%%%%%%%%%%%%%%%%%%%%%%%%%%%%%%%%%
\hypersetup{pdftitle={\ttitle}}
\hypersetup{pdfsubject=\ttitleEn}
\hypersetup{pdfauthor={\tauthor}}
\hypersetup{pdfkeywords=\tkeywordsEn}

%%%%%%%%%%%%%%%%%%%%%%%%%%%%%%%%%%%%%%%%
% postavitev strani
%%%%%%%%%%%%%%%%%%%%%%%%%%%%%%%%%%%%%%%%  

\addtolength{\marginparwidth}{-20pt} % robovi za tisk
\addtolength{\oddsidemargin}{40pt}
\addtolength{\evensidemargin}{-40pt}

\renewcommand{\baselinestretch}{1.3} % ustrezen razmik med vrsticami
\setlength{\headheight}{15pt}        % potreben prostor na vrhu
\renewcommand{\chaptermark}[1]%
{\markboth{\MakeUppercase{\thechapter.\ #1}}{}} \renewcommand{\sectionmark}[1]%
{\markright{\MakeUppercase{\thesection.\ #1}}} \renewcommand{\headrulewidth}{0.5pt} \renewcommand{\footrulewidth}{0pt}
\fancyhf{}
\fancyhead[LE,RO]{\sl \thepage} 
%\fancyhead[LO]{\sl \rightmark} \fancyhead[RE]{\sl \leftmark}
\fancyhead[RE]{\sc \tauthor}              % dodal Solina
\fancyhead[LO]{\sc Diplomska naloga}     % dodal Solina


\newcommand{\BibLaTeX}{{\sc Bib}\LaTeX}
\newcommand{\BibTeX}{{\sc Bib}\TeX}

%%%%%%%%%%%%%%%%%%%%%%%%%%%%%%%%%%%%%%%%
% naslovi
%%%%%%%%%%%%%%%%%%%%%%%%%%%%%%%%%%%%%%%%  

\newcommand{\autfont}{\Large}
\newcommand{\titfont}{\LARGE\bf}
\newcommand{\clearemptydoublepage}{\newpage{\pagestyle{empty}\cleardoublepage}}
\setcounter{tocdepth}{1}	      % globina kazala

%%%%%%%%%%%%%%%%%%%%%%%%%%%%%%%%%%%%%%%%
% konstrukti
%%%%%%%%%%%%%%%%%%%%%%%%%%%%%%%%%%%%%%%%  
\newtheorem{izrek}{Izrek}[chapter]
\newtheorem{trditev}{Trditev}[izrek]
\newenvironment{dokaz}{\emph{Dokaz.}\ }{\hspace{\fill}{$\Box$}}


%%%%%%%%%%%%%%%%%%%%%%%%%%%%%%%%%%%%%%%%%%%%%%%%%%%%%%%%%%%%%%%%%%%%%%%%%%%%%%%
%% PDF-A
%%%%%%%%%%%%%%%%%%%%%%%%%%%%%%%%%%%%%%%%%%%%%%%%%%%%%%%%%%%%%%%%%%%%%%%%%%%%%%%

%%%%%%%%%%%%%%%%%%%%%%%%%%%%%%%%%%%%%%%% 
% define medatata
%%%%%%%%%%%%%%%%%%%%%%%%%%%%%%%%%%%%%%%% 
\def\Title{\ttitle}
\def\Author{\tauthor, lb4129@student.uni-lj.si}
\def\Subject{\ttitleEn}
\def\Keywords{\tkeywordsEn}

%%%%%%%%%%%%%%%%%%%%%%%%%%%%%%%%%%%%%%%% 
% \convertDate converts D:20080419103507+02'00' to 2008-04-19T10:35:07+02:00
%%%%%%%%%%%%%%%%%%%%%%%%%%%%%%%%%%%%%%%% 
\def\convertDate{%
    \getYear
}

{\catcode`\D=12
 \gdef\getYear D:#1#2#3#4{\edef\xYear{#1#2#3#4}\getMonth}
}
\def\getMonth#1#2{\edef\xMonth{#1#2}\getDay}
\def\getDay#1#2{\edef\xDay{#1#2}\getHour}
\def\getHour#1#2{\edef\xHour{#1#2}\getMin}
\def\getMin#1#2{\edef\xMin{#1#2}\getSec}
\def\getSec#1#2{\edef\xSec{#1#2}\getTZh}
\def\getTZh +#1#2{\edef\xTZh{#1#2}\getTZm}
\def\getTZm '#1#2'{%
    \edef\xTZm{#1#2}%
    \edef\convDate{\xYear-\xMonth-\xDay T\xHour:\xMin:\xSec+\xTZh:\xTZm}%
}

%\expandafter\convertDate\pdfcreationdate 

%%%%%%%%%%%%%%%%%%%%%%%%%%%%%%%%%%%%%%%%
% get pdftex version string
%%%%%%%%%%%%%%%%%%%%%%%%%%%%%%%%%%%%%%%% 
\newcount\countA
\countA=\pdftexversion
\advance \countA by -100
\def\pdftexVersionStr{pdfTeX-1.\the\countA.\pdftexrevision}


%%%%%%%%%%%%%%%%%%%%%%%%%%%%%%%%%%%%%%%%
% XMP data
%%%%%%%%%%%%%%%%%%%%%%%%%%%%%%%%%%%%%%%%  
\usepackage{xmpincl}
%\includexmp{pdfa-1b}

%%%%%%%%%%%%%%%%%%%%%%%%%%%%%%%%%%%%%%%%
% pdfInfo
%%%%%%%%%%%%%%%%%%%%%%%%%%%%%%%%%%%%%%%%  
\pdfinfo{%
    /Title    (\ttitle)
    /Author   (\tauthor, lb4129@student.uni-lj.si)
    /Subject  (\ttitleEn)
    /Keywords (\tkeywordsEn)
    /ModDate  (\pdfcreationdate)
    /Trapped  /False
}

%%%%%%%%%%%%%%%%%%%%%%%%%%%%%%%%%%%%%%%%
% znaki za copyright stran
%%%%%%%%%%%%%%%%%%%%%%%%%%%%%%%%%%%%%%%%  

\newcommand{\CcImageCc}[1]{%
	\includegraphics[scale=#1]{cc_cc_30.pdf}%
}
\newcommand{\CcImageBy}[1]{%
	\includegraphics[scale=#1]{cc_by_30.pdf}%
}
\newcommand{\CcImageSa}[1]{%
	\includegraphics[scale=#1]{cc_sa_30.pdf}%
}

%%%%%%%%%%%%%%%%%%%%%%%%%%%%%%%%%%%%%%%%%%%%%%%%%%%%%%%%%%%%%%%%%%%%%%%%%%%%%%%
%%%%%%%%%%%%%%%%%%%%%%%%%%%%%%%%%%%%%%%%%%%%%%%%%%%%%%%%%%%%%%%%%%%%%%%%%%%%%%%


\begin{document}
\selectlanguage{slovene}
\frontmatter
\setcounter{page}{1} %
\renewcommand{\thepage}{}       % preprečimo težave s številkami strani v kazalu

%%%%%%%%%%%%%%%%%%%%%%%%%%%%%%%%%%%%%%%%
%naslovnica
 \thispagestyle{empty}%
   \begin{center}
    {\large\sc Univerza v Ljubljani\\%
%      Fakulteta za elektrotehniko\\% za študijski program Multimedija
%      Fakulteta za upravo\\% za študijski program Upravna informatika
      Fakulteta za računalništvo in informatiko\\%
%      Fakulteta za matematiko in fiziko\\% za študijski program Računalništvo in matematika
     }
    \vskip 10em%
    {\autfont \tauthor\par}%
    {\titfont \ttitle \par}%
    {\vskip 3em \textsc{DIPLOMSKO DELO\\[5mm]         % dodal Solina za ostale študijske programe
%    VISOKOŠOLSKI STROKOVNI ŠTUDIJSKI PROGRAM\\ PRVE STOPNJE\\ RAČUNALNIŠTVO IN INFORMATIKA}\par}%
     VISOKOŠOLSKI  ŠTUDIJSKI PROGRAM\\ PRVE STOPNJE\\ RAČUNALNIŠTVO IN INFORMATIKA}\par}%
%    INTERDISCIPLINARNI UNIVERZITETNI\\ ŠTUDIJSKI PROGRAM PRVE STOPNJE\\ MULTIMEDIJA}\par}%
%    INTERDISCIPLINARNI UNIVERZITETNI\\ ŠTUDIJSKI PROGRAM PRVE STOPNJE\\ UPRAVNA INFORMATIKA}\par}%
%    INTERDISCIPLINARNI UNIVERZITETNI\\ ŠTUDIJSKI PROGRAM PRVE STOPNJE\\ RAČUNALNIŠTVO IN MATEMATIKA}\par}%
    \vfill\null%
% izberite pravi habilitacijski naziv mentorja!
    {\large \textsc{Mentor}: doc. dr. Jurij Mihelič\par}%
   %{\large \textsc{Somentor}:  viš. pred./doc./izr. prof./prof. dr.  Martin Krpan \par}%
    {\vskip 2em \large Ljubljana, \the\year \par}%
\end{center}
% prazna stran
%\clearemptydoublepage      
% izjava o licencah itd. se izpiše na hrbtni strani naslovnice

%%%%%%%%%%%%%%%%%%%%%%%%%%%%%%%%%%%%%%%%
%copyright stran
%%%%%%%%%%%%%%%%%%%%%%%%%%%%%%%%%%%%%%%%
\newpage
\thispagestyle{empty}

\vspace*{5cm}
{\small \noindent
To delo je ponujeno pod licenco \textit{Creative Commons Priznanje avtorstva-Deljenje pod enakimi pogoji 2.5 Slovenija} (ali novej\v so razli\v cico).
To pomeni, da se tako besedilo, slike, grafi in druge sestavine dela kot tudi rezultati diplomskega dela lahko prosto distribuirajo,
reproducirajo, uporabljajo, priobčujejo javnosti in predelujejo, pod pogojem, da se jasno in vidno navede avtorja in naslov tega
dela in da se v primeru spremembe, preoblikovanja ali uporabe tega dela v svojem delu, lahko distribuira predelava le pod
licenco, ki je enaka tej.
Podrobnosti licence so dostopne na spletni strani \href{http://creativecommons.si}{creativecommons.si} ali na Inštitutu za
intelektualno lastnino, Streliška 1, 1000 Ljubljana.

\vspace*{1cm}
\begin{center}% 0.66 / 0.89 = 0.741573033707865
\CcImageCc{0.741573033707865}\hspace*{1ex}\CcImageBy{1}\hspace*{1ex}\CcImageSa{1}%
\end{center}
}

\vspace*{1cm}
{\small \noindent
Izvorna koda diplomskega dela, njeni rezultati in v ta namen razvita programska oprema je ponujena pod licenco GNU General Public License,
različica 3 (ali novejša). To pomeni, da se lahko prosto distribuira in/ali predeluje pod njenimi pogoji.
Podrobnosti licence so dostopne na spletni strani \url{http://www.gnu.org/licenses/}.
}

\vfill
\begin{center} 
\ \\ \vfill
{\em
Besedilo je oblikovano z urejevalnikom besedil \LaTeX.}
\end{center}

% prazna stran
\clearemptydoublepage

%%%%%%%%%%%%%%%%%%%%%%%%%%%%%%%%%%%%%%%%
% stran 3 med uvodnimi listi
\thispagestyle{empty}
\
\vfill

\bigskip
\noindent\textbf{Kandidat:} Luka Bajić\\
\noindent\textbf{Naslov:} Inženiring vzporednih algoritmov\\
% vstavite ustrezen naziv študijskega programa!
\noindent\textbf{Vrsta naloge:} Diplomska naloga na visokošolskem programu prve stopnje Računalništvo in informatika \\
% izberite pravi habilitacijski naziv mentorja!
\noindent\textbf{Mentor:} doc. dr. Jurij Mihelič\\
%\noindent\textbf{Somentor:} isto kot za mentorja

\bigskip
\noindent\textbf{Opis:}\\
V diplomski nalogi \cite{diplomaPremk} je študent preizkusil različne tehnike za povzporejanje algoritma za izračun DTW (dynamic time warping) razdalje, ki sem jih nato še dopolnil. Predvidevam, da bi enake tehnike bile uporabne tudi za nekatere druge probleme, kot je razdalja LCS (longest common subsequence) ali Levenshteinova razdalja. V nalogi je torej treba te tehnike sprogramirati za omenjena problema (lahko le en problem, če bo veliko dela), nato pa jih smiselno eksperimentalno ovrednotiti.

\bigskip
\noindent\textbf{Title:} Parallel algorithm engineering

\bigskip
\noindent\textbf{Description:}\\
opis diplome v angleščini

\vfill



\vspace{2cm}

% prazna stran
\clearemptydoublepage

% zahvala
\thispagestyle{empty}\mbox{}\vfill\null\it%
\noindent
Zahvaljujem se mentorju doc. dr. Juriju Miheliču za pomoč pri izdelavi diplomske naloge in Urošu Koritniku za moralno podporo in koristne nasvete tekom študija. 
\rm\normalfont

% prazna stran
\clearemptydoublepage

%%%%%%%%%%%%%%%%%%%%%%%%%%%%%%%%%%%%%%%%
% posvetilo, če sama zahvala ne zadošča :-)
\thispagestyle{empty}\mbox{}{\vskip0.20\textheight}\mbox{}\hfill\begin{minipage}{0.55\textwidth}%
%Svoji dragi Alenčici.
\normalfont\end{minipage}

% prazna stran
\clearemptydoublepage


%%%%%%%%%%%%%%%%%%%%%%%%%%%%%%%%%%%%%%%%
% kazalo
\pagestyle{empty}
\def\thepage{}% preprečimo težave s številkami strani v kazalu
\tableofcontents{}


% prazna stran
\clearemptydoublepage

%%%%%%%%%%%%%%%%%%%%%%%%%%%%%%%%%%%%%%%%
% seznam kratic

\chapter*{Seznam uporabljenih kratic}

\noindent\begin{tabular}{p{0.11\textwidth}|p{.39\textwidth}|p{.39\textwidth}}    % po potrebi razširi prvo kolono tabele na račun drugih dveh!
  {\bf kratica} & {\bf angleško}                              & {\bf slovensko} \\ \hline
  {\bf LCS}      & longest common subsequence               & najdaljša skupna podsekvenca \\
  {\bf JSON} & JavaScript Object Notation & JavaScript objektna notacija \\
  {\bf SVM}   & support vector machine              & metoda podpornih vektorjev \\
%  \dots & \dots & \dots \\
\end{tabular}


% prazna stran
\clearemptydoublepage

%%%%%%%%%%%%%%%%%%%%%%%%%%%%%%%%%%%%%%%%
% povzetek
\phantomsection
\addcontentsline{toc}{chapter}{Povzetek}
\chapter*{Povzetek}

\noindent\textbf{Naslov:} \ttitle
\bigskip

\noindent\textbf{Avtor:} \tauthor
\bigskip

%\noindent\textbf{Povzetek:} 
\noindent Algoritmi za iskanje podobnosti med nizi so pomemben gradnik raznovrstnih aplikacij, ki so dandanes v vsakodnevni uporabi, od preverjevalnikov črkovanja do orodij za kontrolo verzij (npr. Git). Diplomsko delo opiše različne pristope optimizacije dveh tovrstnih algoritmov: Levenshteinove razdalje in najdaljše skupne podsekvence (??), z uporabo koncepta večnitnosti in dveh temelnjih pristopov paralelizacije - diagonalni in naprej-nazaj. Glavni cilj je časovna pohitritev algoritmov, ki jo ustrezno izmerimo in primerjamo nove implementacije z obstoječimi. Za pristop naprej-nazaj razvijemo tudi prostorsko optimizirano metodo, ki je smiselna za uporabo pri ogromnih količinah vhodnih podatkov. 

\bigskip

\noindent\textbf{Ključne besede:} \tkeywords.
% prazna stran
\clearemptydoublepage

%%%%%%%%%%%%%%%%%%%%%%%%%%%%%%%%%%%%%%%%
% abstract
\phantomsection
\selectlanguage{english}
\addcontentsline{toc}{chapter}{Abstract}
\chapter*{Abstract}

\noindent\textbf{Title:} \ttitleEn
\bigskip

\noindent\textbf{Author:} \tauthor
\bigskip

%\noindent\textbf{Abstract:} 
\noindent This sample document presents an approach to typesetting your BSc thesis using \LaTeX. 
A proper abstract should contain around 100 words which makes this one way too short.
\bigskip

\noindent\textbf{Keywords:} \tkeywordsEn.
\selectlanguage{slovene}
% prazna stran
\clearemptydoublepage

%%%%%%%%%%%%%%%%%%%%%%%%%%%%%%%%%%%%%%%%
\mainmatter
\setcounter{page}{1}
\pagestyle{fancy}

\chapter{Uvod}

V diplomskem delu so opisani različni pristopi za optimizacijo algoritmov za računanje Levenshteinove razdalje in najdaljše skupne podsekvence (??) (LCS - longest common subsequence) z uporabo vzporednega procesiranja. 

\section{Motivacija}



\section{Struktura}

Drugo poglavje opiše algoritma, njuno praktično uporabnost in njuno implementacijo z uporabo dinamičnega programiranja. Tretje poglavje predstavi dva glavna pristopa paralelizacije - diagonalni in naprej-nazaj. V četrtem poglavju so prikazani rezultati meritev na različnih velikostih vhodnih podatkov. V petem poglavju sledijo še sklepi in ideje za nadalnji razvoj. 

\chapter{Metode in paradigme}

V tem poglavju je predstavljena implementacija algoritma za računanje Levenshteinove razdalje in LCS razdalje z uporabo dinamičnega programiranja. V tem kontekstu je namenjena predvsem za časovno primerjavo med sekvenčnim in vzporednim algoritmom. 

\section{Dinamično programiranje}



\section{LCS}

\section{Levenshteinova razdalja}

\subsection{Primer uporabe}

Primer praktične uporabe \cite{documentsimilarity}


\chapter{Vzporedni algoritmi}

Osnovna struktura vzporednega algoritma za oba problema je enaka kot njuna rešitev iz prejšnjega poglavja - še vedno uporabljamo dvodimenzionalno tabelo za hranjenje vmesnih vrednosti, torej je prostorska zahtevnost enaka - O(n*m). Razlika je samo v tem, da vrednosti, ki niso medsebojno odvisne, lahko računamo vzporedno z uporabo koncepta večnitnosti (angl. multithreading). Za to obstajata dva glavna pristopa - diagonalni in naprej-nazaj. 

Vsega skupaj smo implementirali 14 različic algoritmov:
\begin{itemize}
  \item LCS naprej (dinamično programiranje)
  \item LCS naprej s prostorsko optimizacijo
  \item LCS nazaj (dinamično programiranje)
  \item LCS nazaj s prostorsko optimizacijo
  \item LCS naprej-nazaj (2 niti)
  \item LCS naprej-nazaj s prostorsko optimizacijo (2 niti)
  \item LCS diagonalno (poljubno število niti)
  \item Levenshteinova razdalja naprej (dinamično programiranje)
  \item Levenshteinova razdalja naprej s prostorsko optimizacijo
  \item Levenshteinova razdalja nazaj (dinamično programiranje)
  \item Levenshteinova razdalja nazaj s prostorsko optimizacijo
  \item Levenshteinova razdalja naprej-nazaj (2 niti)
  \item Levenshteinova razdalja naprej-nazaj s prostorsko optimizacijo (2 niti)
  \item Levenshteinova razdalja diagonalno (poljubno število niti)
\end{itemize}

\section{Pristop naprej-nazaj}

V prejšnjem poglavju smo po tabeli iterirali od leve proti desni, ter od vrha navzdol, torej za i = 0, 1, ..., n in j = 0, 1, ..., m, nakar smo prišli do končne rešitve na poziciji (i, j). Brez težav lahko počnemo ravno obratno - iteriramo od i = n, n-1, ... 0 in j = m, m-1, ..., 0 in v tem primeru dobimo rešitev na poziciji (0, 0). Izkaže se, da lahko z uporabo dveh niti poženemo oba načina istočasno. Tabelo v abstraktnem smislu prepolovimo na dva dela - ena nit računa zgornjo polovico, druga pa spodnjo. Na ta način, vsaj teoretično, dosežemo dvakratno pohitritev. Na koncu moramo še združiti zadnjo vrstico prve niti z zadnjo vrstico druge niti z združevalni funkcijo [koda], ki pa je časovne zahtevnosti O(n), ker gre dejansko samo za dve enodimenzionalni tabeli, tako da nima bistvenega vpliva na čas izvajanja celotnega algoritma. \ref{LCSnaprej}

\begin{lstlisting}[language=C++, caption={Algoritem LCS naprej}, captionpos=b, label=LCSnaprej]
int forward_LCS (string str1, string str2, int row, int column) {
    for(int i = 0; i < row; i++) 
        for (int j = 0; j < column; j++) 
            if (i == 0 || j == 0)
                arr[i][j] = 0;
            else if (str1[i-1] == str2 [j-1])
                arr[i][j] = 1 + arr[i-1][j-1];
            else
                arr[i][j] = max(arr[i-1][j], arr[i][j-1]);

    return arr[row-1][column-1];
}
\end{lstlisting}






\section{Prostorska optimizacija}

Pri delu z ogromno količino vhodnih podatkov se lahko zgodi da velikost tabele presega kapaciteto pomnilnika. Že za 2 niza dolžine 20000 se izkaže, da če želimo hraniti celotno dvodimenzionalno tabelo v pomnilniku, mora biti velikost le-tega približno 16GB, če predpostavimo da spremenljivke podatkovnega tipa int zasedejo 4 bajte na danem sistemu. Pri naših meritvah, ki so podrobneje opisane v poglavju 4, smo se temu problemu izognili z uporabo spremenljivk podatkovnega tipa unsigned short int, ki zasedajo samo 2 bajta v pomnilniku. Pri tem smo morali biti pozorni, da ne pride do overflowa (??), saj lahko z dvema bajtoma hranimo maksimalno vrednost 65535. Vendar ker vemo, da je tako za LCS kot za Levenshteinovo razdaljo zgornja meja rezultata manjša ali enaka dolžini daljšega izmed vhodnih nizov, je ta pristop vseeno smiselen za vhodne nize krajše od dolžine 65536 (čeprav nizi te dolžine kljub vsemu presegajo kapaciteto pomnilnika na povprečnem osebnem računalniku, na katerem smo poganjali teste). 

Vseeno pa je smiselno najti splošnejšo rešitev. Izkaže se, da pri pristopu naprej-nazaj, oziroma tudi brez uporabe večnitnosti pri pristopih naprej in nazaj, lahko bistveno zmanjšamo porabo prostora tako, da v pomnilniku hranimo samo vrstico, ki jo trenutno računamo in prejšnjo vrstico, ostale so nepotrebne in jih lahko zavržemo. Kot je razvidno iz rezultatov v 4. poglavju, smo s tem nekoliko poslabšali časovno zahtevnost, saj moramo v vsaki iteraciji kreirati dodatno enodimenzionalno tabelo, oziroma prepisovati vrednosti. Vsekakor pa je ta način zelo uporaben za ogromne vhodne podatke. 

\section{Diagonalni pristop}

Pomanjkljivost pristopa naprej-nazaj je, da dejansko uporablja samo dve niti, oziroma ga je težko dodatno paralelizirati, ker se morajo niti v tem primeru med seboj čakati in ne dosežemo maksimalne optimalnosti, oziroma se čas morda celo poslabša zaradi prevelikega overheada. 

Zato če želimo uporabiti več kot dve niti, raje koristimo diagonalni pristop, ki nam teoretično omogoča hkratno uporabo do k niti (k = dolžina diagonale = dolžina krajšega izmed dveh nizov). Ta pristop deluje, ker se izkaže, da so elementi na diagonalah vedno medsebojno neodvisni, algoritem namreč dostopa samo do podatkov v celicah (i, j-1), (i-1, j) in (i-1, j-1), paziti moramo samo, da se diagonale računajo v pravem vrstnem redu, torej da začnemo s celico (0, 0) in končamo s celico (n, m). To najlažje dosežemo tako, da z zunanjo zanko iteriramo po diagonalah (teh je vedno n+m-1) in nato v vsaki iteraciji vzporedno računamo vrednosti vseh celic na posamezni diagonali. 

[slikovni prikaz]











\chapter{Rezultati}

V tem poglavju so prikazane meritve časov izvajanja za posamezne algoritme in medsebojna primerjava le-teh. Glede na to, da je večnitenje zelo odvisno od same arhitekture sistema na katerem izvajamo programsko kodo, pa tudi od raznih zunanjih dejavnikov (npr. scheduling operacijskega sistema), je predvsem pri tabelah manjših velikosti smiselno, da posamezno funkcijo kličemo večkrat in kot končni rezultat vzamemo povprečje vseh meritev. 

\section{Generiranje podatkov}

Vhodni podatki so enakega tipa tako za LCS kot za Levenshteinovo razdaljo: teoretično gre za poljubno število znakovnih nizov, vendar smo se za potrebe tega diplomskega dela omejili na delo z točno dvema nizoma, ki vsebujeta znake iz angleške abecede. Na same algoritme ne bi imelo nobenega vpliva, če bi uporabili celotno ASCII tabelo ali katerikoli drug nabor znakov, edina razlika bi bila v tem, da bi bile izračunane razdalje v povprečju krajše zaradi manjše teoretične verjetnosti da se posamezen znak pojavi v obeh nizih. Posledično bi se izvedlo manj klicev funkcij min oz. max, vendar je kar se tiče časovne zahtevnosti ta pohitritev zanemarljiva. 

\section{Delovno okolje}

VSCode
G++
-O3
AMD
16GB

\section{Teoretične pohitritve/Časovne zahtevnosti}



\section{Primerjava sekvenčnih algoritmov}



\section{Primerjava vzporednih algoritmov}



\section{Celotna primerjava}






\chapter{Zaključki}

\section{Sklepi}

\section{Nadalnji razvoj}

Kot je razvidno iz poglavja 4., smo z vzporednimi algoritmi dosegli občutno pohitritev že na običajnih osebnih računalnikih. V nadalnje bi bilo zanimivo preveriti ali je čas izvajanja morda še krajši na zmogljivejših sistemih z ogromno kolilčino jeder, oziroma ali je mogoče algoritme še dodatno prilagoditi za tovrstne sisteme. 


%\cleardoublepage
%\addcontentsline{toc}{chapter}{Literatura}

% če imaš težave poravnati desni rob bibliografije, potem odkomentiraj spodnjo vrstico
\raggedright

%\printbibliography[heading=bibintoc,type=article,title={Članki v revijah}]

%\printbibliography[heading=bibintoc,type=inproceedings,title={Članki v zbornikih}]

%\printbibliography[heading=bibintoc,type=incollection,title={Poglavja v knjigah}]

% v zadnji verziji diplomskega dela običajno združiš vse tri vrste referenc v en sam seznam in
% izpustiš delne sezname
\printbibliography[heading=bibintoc,title={Literatura}]

\end{document}

